\originalchapter{8.04 作业5：不确定性与势阱}
来源为8.04 Spring 2016 Assignment 5, 页面与标题为2016春, PDF发布行印为2015年3月11日, 截止行是2016年3月18日中午; 保留该日期差异, 不默改原件. 全部八题共65分, 以下为\textbf{AI独立解答}.
\originalsection{第1题：实Gaussian（5分）}
\begin{exercise}
令$\psi=N\e^{-x^2/(2a^2)}$, $a>0$. (a)在位置表象计算$\Delta x,\Delta p$, 检验不确定性是否饱和. (b)求动量波函数, 验证归一化并重新求$\Delta p$.
\end{exercise}
\begin{solution}
$N=(\pi a^2)^{-1/4}$, 偶性给两均值0. 由$I(b)=\int\e^{-bx^2}\dd x=\sqrt{\pi/b}$,$-I\prime(b)=\int x^2\e^{-bx^2}\dd x=\sqrt\pi/(2b^{3/2})$, 取$b=1/a^2$并除归一积分, 得$\langle x^2\rangle=a^2/2$. 用$\langle p^2\rangle=\hbar^2\int|\psi'|^2\dd x=\hbar^2/(2a^2)$, 得$\Delta x=a/\sqrt2,\Delta p=\hbar/(\sqrt2a)$, 乘积$\hbar/2$. Fourier积分给$\phi(p)=\sqrt{a/(\sqrt\pi\hbar)}\e^{-a^2p^2/(2\hbar^2)}$. 积分$|\phi|^2$为1, 其二阶矩为$\hbar^2/(2a^2)$, 与位置计算一致.
\end{solution}
\originalsection{第2题：复Gaussian与相关性（10分）}
\begin{exercise}
把$a$替换为复$\Delta$, $\operatorname{Re}(\Delta^2)>0$, $N$取实正. (a)求宽度, 以$|\Delta|$和$\operatorname{Re}(\Delta^2)$表示. (b)Fourier变换, 检验Parseval归一化, 并在动量表象重算$\Delta p$. (c)令$\Delta=|\Delta|\e^{\ii\phi_\Delta}$, 求不确定性乘积并讨论$\phi_\Delta=0,\pi/4$. (d)用于作业4自由Gaussian, 验证动量宽度不随时间变.
\end{exercise}
\begin{solution}
设$D=\Delta^2$, $A=\operatorname{Re}(1/D)=\operatorname{Re}D/|D|^2>0$, $N=(A/\pi)^{1/4}$. 共轭分式$1/z=z^*/|z|^2$证明题给恒等式. $\langle x^2\rangle=1/(2A)$, 而$\psi'=-x\psi/D$, 故
\[
\Delta x=\frac{|\Delta|^2}{\sqrt{2\operatorname{Re}(\Delta^2)}},\qquad
\Delta p=\frac{\hbar}{\sqrt{2\operatorname{Re}(\Delta^2)}},\qquad
\phi(p)=\frac{N\sqrt D}{\sqrt\hbar}\e^{-\Delta^2p^2/(2\hbar^2)}.
\]
固定Fourier变换取$\operatorname{Re}\sqrt D>0$的平方根; $D=\Delta^2$允许原参数$\Delta$正负两代表, 不能未经选择直接把根写成$\Delta$. 这个选择使变换等号精确, 而不仅到整体相位. 模方动量密度为$N^2|\Delta|^2\e^{-\operatorname{Re}D\,p^2/\hbar^2}/\hbar$, 用前题$-I\prime(b)/I(b)=1/(2b)$并取$b=\operatorname{Re}D/\hbar^2$, 积分为1（亦为Parseval检验）, 两均值由偶性为0, 动量方差$\hbar^2/[2\operatorname{Re}D]$与位置表象计算相同. 乘积为$\hbar/[2\cos(2\phi_\Delta)]$（要求分母正）. $\phi_\Delta=0$取等号; $\phi_\Delta\to\pi/4$时位置宽度发散, 端点本身不可归一, 不能当成有限宽度态. 作业4有$D(t)=2a^2(1+\ii t/\tau)$, 实部恒为$2a^2$, 所以$\Delta p=\hbar/(2a)$始终不变. 位置展宽来自相关相位, 不来自动量分布变宽.
\end{solution}
\originalsection{第3题：盒中定态的宽度（15分）}
\begin{exercise}
区间$[0,a]$有无限高墙, $\Psi_n=N\sin(n\pi x/a)\e^{-\ii\phi_n(t)}$. (a)求相位和归一化. (b)求位置一、二阶矩及宽度. (c)求$\langle p\rangle,\langle p^2\rangle$及宽度,$n=1,2,\ldots$. (d)判断不确定性是否满足及饱和. (e)说明上述量是否随时间变.
\end{exercise}
\begin{solution}
$N=\sqrt{2/a}$, $E_n=\hbar^2n^2\pi^2/(2ma^2)$, $\phi_n(t)=E_nt/\hbar+\phi_n(0)$. 用$\sin^2u=(1-\cos2u)/2$. 令$b=2n\pi/a$, 两次分部积分给$\int_0^a x^2\cos bx\dd x=2a/b^2=a^3/(2n^2\pi^2)$; 一阶余弦矩为0. 将这两项代回归一密度积分得
\[
\langle x\rangle=\frac a2,\quad\langle x^2\rangle=a^2\left(\frac13-\frac1{2n^2\pi^2}\right),\quad
\Delta x=a\sqrt{\frac1{12}-\frac1{2n^2\pi^2}}.
\]
$\int\psi\psi'\dd x=0$给$\langle p\rangle=0$, $-\hbar^2\psi\prime\prime=(n\pi\hbar/a)^2\psi$给$\langle p^2\rangle=(n\pi\hbar/a)^2$, 因而$\Delta p=n\pi\hbar/a$. 乘积$\hbar\sqrt{n^2\pi^2/12-1/2}$, 最小$n=1$约$0.568\hbar>\hbar/2$, 不饱和. 定态仅乘整体相位, 全部矩不随时间变. 这里的动量矩可由全线延零波函数的动量分布解释; 不把盒内零边界的一阶导数算符误称为自伴动量观测量.
\end{solution}
\originalsection{第4、5题：半线硬墙与上台阶（各5、10分）}
\begin{exercise}
第4题: $x\leq0$为硬墙,$x>0$自由, 求不可归一化定态与能量. 第5题: 全线$V=0$（$x\leq0$）、$V_0$（$x>0$）, 求$0<E<V_0$的定态.
\end{exercise}
\begin{solution}
第4题零边界迫使$u_k(x)=A\sin kx$（$x>0$）, $E=\hbar^2k^2/(2m),k>0$. 不存在负能正规解; $E=0$线性解不提供可归一态. 第5题以单位左入射写$u=\e^{\ii kx}+r\e^{-\ii kx}$（左）,$u=t\e^{-\kappa x}$（右）, $k=\sqrt{2mE}/\hbar$, $\kappa=\sqrt{2m(V_0-E)}/\hbar$. 连续条件$1+r=t$, $\ii k(1-r)=-\kappa t$给$r=(k-\ii\kappa)/(k+\ii\kappa)$,$t=2k/(k+\ii\kappa)$. $|r|=1$, 右侧有概率尾巴却无净流. 乘任意复常数得到同一能量散射解, 没有从右衰减支路来的独立入射态.
\end{solution}
\originalsection{第6题：硬墙旁的半势阱（10分）}
\begin{exercise}
$x<0$硬墙,$0<x<a$势为$-V_0$,$x>a$为0. 求$E<0$束缚态, 判断是否总有束缚态; 令$z_0^2=2ma^2V_0/\hbar^2$, 求三束缚态阈值并联系宽$2a$全线势阱.
\end{exercise}
\begin{solution}
令$k=\sqrt{2m(E+V_0)}/\hbar$, $\kappa=\sqrt{-2mE}/\hbar$, 则内部$A\sin kx$, 外部$A\sin(ka)\e^{-\kappa(x-a)}$. 匹配导数得$k\cot ka=-\kappa$. 写$z=ka,w=\kappa a$, 有$z\cot z=-\sqrt{z_0^2-z^2}$, $E=\hbar^2z^2/(2ma^2)-V_0$. 根分别出现在$((n-1/2)\pi,n\pi)$与$z<z_0$的交集. 第一根需$z_0>\pi/2$, 所以并非总有束缚态. 第三个阈值为$5\pi/2$, 严格有三个需$z_0>5\pi/2$且$z_0\leq7\pi/2$（下一态在上端恰为零能、不正规）. 在等号处新态$E=0$的外尾为常数, 不可归一; 因而没有取到的最小值, 只有下确界. 将半线态奇延拓到负轴, 正是宽$2a$有限阱的\textbf{奇}态; 偶态没有对应, 不能把全阱所有束缚态数搬过来.
\end{solution}
\originalsection{第7题：模拟氢能量的方阱（5分）}
\begin{exercise}
全线$|x|<a_0$势$-V_0$, 外为0, 要使基态为$E_0=-e^2/(2a_0)=-13.6$ eV, 求$V_0$.
\end{exercise}
\begin{solution}
氢关系$|E_0|=\hbar^2/(2ma_0^2)$给$\kappa a_0=1$. 偶基态条件$z\tan z=1$, $0<z<\pi/2$, 数值根$z=0.860334$, 故$V_0=|E_0|(1+z^2)\simeq23.67$ eV. 此处只匹配一个能量与长度, 一维方阱没有氢的整个谱或角动量结构.
\end{solution}
\originalsection{第8题：为何不能处处禁阻（5分）}
\begin{exercise}
实定态满足$-\hbar^2\psi''/(2m)+(V-E)\psi=0$. (a)从$\langle H\rangle$证明能量高于势最小值. (b)用波函数曲率说明处处$E<V$不能形成正规束缚态.
\end{exercise}
\begin{solution}
在正规归一化态和消失边界项下,$E=\hbar^2\int|\psi'|^2/(2m)+\int V|\psi|^2\geq V_{\min}$. 全线正规态动能不可能为0, 因为那会要求非零常函数, 不可归一, 所以严格大于$V_{\min}$. 若处处$V-E>0$, 则$\psi''$与$\psi$同号: 正的段向上凸, 负的段向下凸. 一条两端衰减的非零连续曲线必在正最大处有非正曲率, 或在负最小处有非负曲率, 与方程冲突. 这个几何论证与积分正性是同一限制, 不能由禁阻区指数解在单个区间存在推出整个束缚态存在.
\end{solution}
